\documentclass[sigconf,nonacm]{acmart}

\usepackage{amsmath,amssymb,amsfonts}
\usepackage{graphicx}
\usepackage{booktabs}
\usepackage{algorithm}
\usepackage{algpseudocode}
\usepackage{hyperref}
\usepackage{xcolor}
\usepackage{subcaption}

\begin{document}

\title{GeoBind: Hyperbolic-Euclidean Composite Token Learning\\for Semantic Binding in Text-to-Image Generation\\
\large Supplementary Material}

\maketitle

\renewcommand{\thesection}{\Alph{section}}
\renewcommand{\thetable}{\Alph{section}\arabic{table}}
\renewcommand{\thefigure}{\Alph{section}\arabic{figure}}
\renewcommand{\theequation}{\Alph{section}.\arabic{equation}}

% ============================================================
\section{Algorithm}
\label{sec:algorithm}
% ============================================================

Algorithm~\ref{alg:geobind} presents the complete GeoBind inference procedure.

\begin{algorithm}[t]
\caption{GeoBind Inference}
\label{alg:geobind}
\begin{algorithmic}[1]
\Require Prompt $P$, base T2I model $\mathcal{M}$, curvature $c$, fusion weight $\gamma$, steps $T$, token opt steps $T_{\text{tok}}$, attn opt steps $T_{\text{attn}}$
\Ensure Generated image $\mathbf{x}_0$
\State Parse $P$ with SpaCy $\rightarrow$ entities $\{(n_k, a_k)\}_{k=1}^K$
\State Encode $P \rightarrow C = \{c_1, \ldots, c_L\} \in \mathbb{R}^{L \times d}$ via CLIP
\State Normalize: $\tilde{C} \leftarrow C / \|C\|_2$
\State Map to Poincar\'e ball: $C^{\mathcal{H}} \leftarrow \exp_\mathbf{0}^c(\tilde{C})$
\State Add positional encoding via M\"obius addition (Eq.~10)
\For{each entity $k = 1, \ldots, K$}
    \State Euclidean merge: $\hat{c}_k^{\text{E}} \leftarrow \alpha \cdot c_{n_k} + \beta \cdot \sum_{j} c_{a_{k,j}}$ \Comment{Eqs.~11, 13}
    \State Hyperbolic merge: $\hat{c}_k^{\mathcal{H}} \leftarrow$ attention-based fusion in $\mathcal{H}$ \Comment{Eqs.~12, 14}
    \State Geometry-aware fusion: $\hat{c}_k \leftarrow \hat{c}_k^{\text{E}} + \gamma \cdot \log_\mathbf{0}^c(\hat{c}_k^{\mathcal{H}})$ \Comment{Eq.~15}
\EndFor
\State Replace original tokens with composite tokens $\{\hat{c}_k\}$
\State Encode simplified prompt $P' \rightarrow C'$ (entity names only)
\State Sample $\mathbf{z}_T \sim \mathcal{N}(0, I)$
\For{$t = T, T-1, \ldots, 1$}
    \If{$t > T - T_{\text{tok}}$} \Comment{Token optimization phase}
        \For{iter $= 1, \ldots, 3$}
            \State Compute $\mathcal{L}_{\text{sem}}$ from noise predictions \Comment{Eq.~16}
            \State Compute $\mathcal{L}_{\text{hyp}}$ from geodesic distances \Comment{Eq.~17}
            \State Update $\hat{c}_k \leftarrow \hat{c}_k - \eta \nabla_{\hat{c}_k}(\lambda_1 \mathcal{L}_{\text{sem}} + \lambda_2 \mathcal{L}_{\text{hyp}})$ \Comment{Eq.~18}
        \EndFor
    \EndIf
    \If{$t > T - T_{\text{attn}}$} \Comment{Attention refinement phase}
        \State Compute $\mathcal{L}_{\text{ent}}$ from attention entropy \Comment{Eq.~19--20}
        \State Update $\mathbf{z}_t \leftarrow \mathbf{z}_t - \alpha_t \nabla_{\mathbf{z}_t} \mathcal{L}_{\text{ent}}$
    \EndIf
    \State Replace EOT tokens with those from $C'$ (ETS)
    \State $\mathbf{z}_{t-1} \leftarrow \text{Denoise}(\mathbf{z}_t, C_{\text{merged}}, t)$
\EndFor
\State $\mathbf{x}_0 \leftarrow \text{VAE\_Decode}(\mathbf{z}_0)$
\end{algorithmic}
\end{algorithm}

% ============================================================
\section{Extended Method Details}
\label{sec:extended_method}
% ============================================================

\subsection{Information Coupling in CLIP Embeddings}

As observed by~\cite{hu2024tome}, the semantic content of CLIP text tokens is \emph{coupled}: individual token embeddings (e.g., $c_{\text{dog}}$) already contain leaked information from their context (e.g., ``sunglasses''), and the end-of-text [EOT] tokens further compound this issue by encapsulating the entire prompt's semantics through masked self-attention. This coupling is a root cause of attribute confusion when multiple objects coexist in the prompt.

Our End Token Substitution (ETS) mechanism directly addresses this coupling by replacing the [EOT] tokens from the full prompt with those encoded from a simplified prompt containing only entity names (e.g., ``a cat and a dog''), effectively removing leaked attribute information from the end-of-text representation.

\subsection{Semantic Additivity Extends to Hyperbolic Space}

Prior work~\cite{hu2024tome} shows that CLIP text embeddings exhibit \emph{semantic additivity}: summing ``dog'' and ``hat'' embeddings yields a token that generates dogs with hats. We extend this insight by replacing Euclidean addition with M\"obius addition (Eq.~4 in the main paper), obtaining composite tokens that combine semantics \emph{and} preserve hierarchical proximity. As demonstrated in Sec.~4.5 of the main paper, this leads to more accurate binding, especially for prompts with deep dependency chains (e.g., ``a dog wearing a red hat with a golden buckle'').

\subsection{Boundary Amplification in High Dimensions}

We provide a formal justification for why hyperbolic geometry helps even in $\mathbb{R}^{2048}$. For two points $\mathbf{x}, \mathbf{y}$ near the boundary of the Poincar\'e ball with $\|\mathbf{x}\|, \|\mathbf{y}\| \to 1/\sqrt{c}$:

\begin{equation}
D_{\text{hyp}}(\mathbf{x}, \mathbf{y}) \approx 2 \log \frac{2\|\mathbf{x} - \mathbf{y}\|}{(1 - c\|\mathbf{x}\|^2)(1 - c\|\mathbf{y}\|^2)},
\label{eq:boundary}
\end{equation}

which diverges logarithmically as the conformal factors $\lambda_c(\mathbf{x}) = 2/(1-c\|\mathbf{x}\|^2)$ grow. This ``boundary amplification'' effect operates regardless of the ambient dimensionality and provides a geometric inductive bias that is absent in Euclidean space. In practice, leaf-node tokens (attributes like ``red hat'') that are assigned to \emph{different} parent entities are pushed exponentially further apart in geodesic distance, even when their Euclidean distances remain similar.

% ============================================================
\section{Implementation Details}
\label{sec:impl_details}
% ============================================================

\setcounter{table}{0}

We provide the complete implementation details for reproducibility.

\paragraph{Syntactic Parsing.} We use SpaCy~\cite{honnibal2020spacy} (model: \texttt{en\_core\_web\_trf}) for dependency parsing to extract entity--attribute pairs. The parser identifies noun chunks and their modifiers (amod, nmod, compound, npadvmod, advmod, acomp) to construct the set of entities $\{(n_k, a_k)\}_{k=1}^K$.

\paragraph{Poincar\'e Ball Parameters.} The curvature is fixed at $c=1.0$. Before the exponential map, all CLIP embeddings are normalized to unit norm ($\tilde{C} = C / \|C\|_2$) to prevent tanh saturation---since raw CLIP embeddings have norms in [15, 50], direct application would cause numerical overflow.

\paragraph{Multi-Head Attention.} We use 8 attention heads with $d=2048$. The projection matrices are fixed as identity matrices ($W_Q = W_K = W_V = I$), making this entirely parameter-free.

\paragraph{Fusion Parameters.} The aggregation weights are $\alpha = 1.1$ (noun), $\beta = 1.2$ (attribute). The hyperbolic--Euclidean fusion weight is $\gamma = 0.1$.

\paragraph{Denoising Configuration.} We run 50 denoising steps with classifier-free guidance scale 7.5 (SDXL) or 3.5 (FLUX.1). Composite token optimization runs for the first 5 steps with 3 gradient iterations per step (learning rate $\eta = 10000$). Attention refinement applies $[4, 4]$ iterations during the first 5 steps.

\paragraph{Loss Weights.} Semantic binding loss weight $\lambda_1 = 1.0$, hyperbolic contrastive loss weight $\lambda_2 = 10^{-6}$, contrastive temperature $\tau = 0.07$.

\paragraph{Evaluation Seeds.} All experiments use 5 random seeds (42, 123, 456, 789, 1024). We report mean $\pm$ standard deviation.

\paragraph{Hardware.} Experiments are conducted on NVIDIA A100 (80 GB) GPUs. Single-image generation takes 12.8--16.4 seconds depending on entity count (Table~7 in the main paper).

% ============================================================
\section{Hyperparameter Sensitivity}
\label{sec:hyperparam}
% ============================================================

\setcounter{table}{0}

We systematically evaluate the sensitivity of four key hyperparameters.

\paragraph{Fusion weight $\gamma$.} We sweep $\gamma \in \{0.01, 0.05, 0.1, 0.2, 0.5, 1.0\}$ (Table~\ref{tab:gamma}). Performance peaks at $\gamma=0.1$ for color and texture, and at $\gamma=0.2$ for shape. Beyond $\gamma=0.5$, performance degrades as the hyperbolic correction overwhelms the Euclidean component, causing numerical instability in cross-attention computation. We set $\gamma=0.1$ as a balanced default. Notably, the \emph{effective} contribution of the hyperbolic branch is larger than $\gamma$ suggests: measuring the $\ell_2$ norm ratio $\|\gamma \cdot \log_0^c(\hat{c}_k^{\text{H}})\| / \|\hat{c}_k^{\text{E}}\|$ across all tokens and prompts, we find a mean effective ratio of $0.23 \pm 0.08$, because the logarithmic map amplifies norms for tokens near the Poincar\'e boundary.

\begin{table}[t]
\centering
\caption{Sensitivity to fusion weight $\gamma$ (BLIP-VQA $\uparrow$, T2I-CompBench, SDXL).}
\label{tab:gamma}
\begin{tabular}{l cccccc}
\toprule
$\gamma$ & 0.01 & 0.05 & \textbf{0.1} & 0.2 & 0.5 & 1.0 \\
\midrule
Color  & .762 & .778 & \textbf{.791} & .785 & .764 & .721 \\
Texture & .665 & .669 & \textbf{.673} & .671 & .658 & .634 \\
Shape  & .547 & .583 & .615 & \textbf{.622} & .608 & .571 \\
\bottomrule
\end{tabular}
\end{table}

\paragraph{Curvature $c$.} We test $c \in \{0.1, 0.5, 1.0, 2.0, 5.0\}$. Performance is stable for $c \in [0.5, 2.0]$, with mild degradation at extreme values. Higher $c$ compresses the Poincar\'e ball, causing numerical precision issues; lower $c$ approaches Euclidean geometry, reducing the hyperbolic benefit. We set $c=1.0$.

\paragraph{Contrastive temperature $\tau$.} We sweep $\tau \in \{0.01, 0.03, 0.07, 0.1, 0.3\}$. Smaller $\tau$ increases the sharpness of the contrastive objective; $\tau=0.07$ provides the best balance. Extremely small $\tau$ ($\leq 0.01$) causes training instability.

\paragraph{Loss weight $\lambda_2$.} We test $\lambda_2 \in \{10^{-8}, 10^{-7}, 10^{-6}, 10^{-5}, 10^{-4}\}$. The optimal range is $[10^{-7}, 10^{-5}]$; larger values cause the contrastive loss to dominate and distort the composite tokens.

% ============================================================
\section{Extended Ablation Analysis}
\label{sec:extended_ablation}
% ============================================================

Building on the ablation study in Table~6 of the main paper, we provide detailed analysis of each component's contribution:

\begin{enumerate}
\item \textbf{Token merging provides a strong foundation.} Config.~B (Euclidean merging only) already improves upon the SDXL baseline across all subsets, confirming the effectiveness of composite tokens~\cite{hu2024tome}.

\item \textbf{Hyperbolic fusion is critical for shape binding.} Comparing E with B, adding the hyperbolic branch improves shape from 0.544 to 0.568 (+2.4\%). Comparing F with D, shape binding jumps from 0.532 to 0.593 (+6.1\%).

\item \textbf{The hyperbolic branch is not a generic perturbation.} Config.~G replaces the hyperbolic branch with a random orthogonal projection of the same dimensionality (generating a correction term $\gamma \cdot R\tilde{C}$ where $R$ is a random orthogonal matrix, $R^\top R = I$). This ``random projection control'' achieves only marginal improvement over ToMe (shape: 0.540 vs.\ 0.532), far below the hyperbolic variant (0.615). This confirms that the gains from hyperbolic geometry are specifically attributable to the hierarchical inductive bias of the Poincar\'e ball.

\item \textbf{Entropy loss drives color binding but can slightly hurt shape.} Adding $\mathcal{L}_{\text{ent}}$ (C vs.\ B) yields the largest single-component gain for color binding (+9.5\%), but shape decreases slightly (0.544$\to$0.531). This trade-off is resolved by our hyperbolic fusion.

\item \textbf{Hyperbolic contrastive loss provides consistent gains.} Comparing the full model with F, $\mathcal{L}_{\text{hyp}}$ improves color (+1.7\%), texture (+0.4\%), and shape (+2.2\%), demonstrating that geodesic separation complements all other components.

\item \textbf{All components are complementary.} The full GeoBind achieves the best performance on every metric. The cumulative gain over ToMe (D) is +3.6/+0.9/+8.3 on color/texture/shape.
\end{enumerate}

% ============================================================
\section{Extended Qualitative Analysis}
\label{sec:extended_qual}
% ============================================================

Following~\cite{rassin2023syngen}, we categorize failure cases of baseline methods into three types:
\begin{itemize}
\item \textbf{Semantic leak in prompt}: an attribute is incorrectly assigned to the wrong entity within the prompt (e.g., ``red'' leaking from ``pear'' to ``apple'').
\item \textbf{Semantic leak out of prompt}: an attribute appears on background elements or unmentioned objects.
\item \textbf{Attribute neglect}: a specified attribute is completely absent from the generated image.
\end{itemize}

Our analysis of 200 GOB-Bench images shows that GeoBind reduces ``semantic leak in prompt'' errors by 73\% compared to the SDXL baseline and by 42\% compared to ToMe. The improvement is most pronounced for prompts with 3+ entities, where the hyperbolic contrastive loss explicitly separates all entity embeddings via pairwise geodesic distances.

% ============================================================
\section{Failure Cases}
\label{sec:failures}
% ============================================================

We identify three main failure modes of GeoBind:

\begin{enumerate}
\item \textbf{SpaCy parsing errors.} Highly informal or grammatically unusual prompts (e.g., ``angry red car next to happy blue truck both racing'') may yield incorrect dependency parses, leading to wrong entity--attribute groupings. This could be mitigated by using LLM-based parsing in future work.

\item \textbf{Extreme entity count.} With $K \geq 6$ entities, the optimization landscape becomes challenging and runtime increases significantly ($>$25\,s). The pairwise contrastive loss scales as $O(K^2)$, making both computation and gradient estimation noisy for large $K$.

\item \textbf{Rare attribute combinations.} Attributes that rarely co-occur in CLIP's training data (e.g., ``pentagonal cat'') cannot be reliably bound regardless of geometric approach, as the base model lacks the visual concept. This is a limitation of the underlying T2I model rather than our binding method.
\end{enumerate}

% ============================================================
\section{Limitations and Future Work}
\label{sec:limitations}
% ============================================================

Our method relies on SpaCy for syntactic parsing, which may fail on informal or non-English text; extending to multilingual settings requires language-specific parsing and multilingual encoders. The curvature $c$ is currently fixed; adaptive curvature selection based on prompt complexity could further improve performance. With $K \geq 6$ entities, optimization becomes challenging. We believe the integration of non-Euclidean geometry into generative models opens a promising direction for improving compositional understanding in T2I systems.

% ============================================================
\section{Additional Related Work: Compositional T2I Benchmarks}
\label{sec:related_benchmarks}
% ============================================================

T2I-CompBench~\cite{huang2023t2icompbench} provides a systematic evaluation suite for compositional generation, covering color, shape, texture, spatial, and non-spatial relationships. DPG-Bench~\cite{hu2024ella} evaluates dense prompt understanding with long descriptions. GenAI-Bench~\cite{li2024genaibench} assesses compositional alignment across diverse categories. ImageReward~\cite{xu2024imagereward} offers a learned human-preference metric. We complement these with our proposed GOB-Bench (Sec.~4.1 of the main paper) to evaluate object binding, a critical dimension not directly addressed by existing benchmarks.

% ============================================================
\section{GOB-Bench Dataset Details}
\label{sec:gobbench}
% ============================================================

\setcounter{table}{0}

GOB-Bench (GeoBind Object Binding Benchmark) contains 200 \emph{object binding} prompts spanning four difficulty levels, designed to systematically evaluate whether each object in a prompt is correctly associated with its sub-objects and accessories (e.g., ``hat'' should belong to ``dog,'' not ``cat'').

\begin{table}[t]
\centering
\caption{GOB-Bench difficulty level statistics.}
\label{tab:gob_stats}
\begin{tabular}{l c c c}
\toprule
Level & \# Objects & \# Sub-objs/Obj & \# Prompts \\
\midrule
Easy     & 2 & 1     & 60 \\
Medium   & 2 & 2     & 60 \\
Hard     & 3 & 1--2  & 40 \\
Complex  & 4+ & 1--2  & 40 \\
\midrule
\textbf{Total} & --- & --- & \textbf{200} \\
\bottomrule
\end{tabular}
\end{table}

\paragraph{Prompt Design Principles.}
Object binding differs from attribute binding (tested by T2I-CompBench) in that each object must be correctly linked to its sub-objects or accessories, rather than adjective modifiers.
\begin{enumerate}
\item \textbf{Template structure}: all prompts follow the ``a [object] with/wearing [item(s)]'' pattern, which creates clear object--sub-object ownership relationships. Example: ``a cat with sunglasses and a dog with a hat.''
\item \textbf{Entity diversity}: objects span animals, people, occupations, and fictional characters to avoid domain bias.
\item \textbf{Sub-object diversity}: sub-objects include accessories (sunglasses, hat, crown, scarf), tools (wand, sword, brush), and wearable items (cape, vest, helmet) that can plausibly belong to the associated object.
\item \textbf{Difficulty calibration}: Easy (2 objects $\times$ 1 item), Medium (2 objects $\times$ 2 items), Hard (3 objects $\times$ 1--2 items), Complex (4+ objects $\times$ 1--2 items). Each level increases the binding ambiguity by adding more competing entities.
\item \textbf{Template-based parsing}: since object binding prompts use ``with/wearing'' (prepositional phrases, not adjective modifiers), we use template-based parsing rather than SpaCy dependency parsing. This ensures deterministic and correct entity--sub-object grouping for all 200 prompts.
\end{enumerate}

\paragraph{Evaluation Protocol.}
Each prompt is evaluated by three independent judges: GPT-4o, Qwen2-VL, and human annotators (3 annotators per image, majority vote). For automated judges, we use the following scoring rubric:
\begin{itemize}
\item \textbf{100}: All objects present with all sub-objects correctly bound to their owners.
\item \textbf{75}: All objects present, but 1 sub-object incorrectly assigned (e.g., hat on cat instead of dog).
\item \textbf{50}: All objects present, but 2+ sub-objects incorrectly assigned.
\item \textbf{25}: One or more objects missing from the image.
\item \textbf{0}: Image does not match the prompt at all.
\end{itemize}
The inter-evaluator agreement (Krippendorff's $\alpha = 0.83$) indicates strong consistency across the three judges.

The complete prompt list is provided in the supplementary file \texttt{gob\_bench\_prompts.txt}.

% ============================================================
\section{Training and Computational Details}
\label{sec:training}
% ============================================================

GeoBind is a \textbf{training-free} method: no model weights are updated. All optimization occurs at inference time on the token embedding space. The key computational aspects are:

\begin{itemize}
\item \textbf{No training data required}: GeoBind operates directly on the pre-trained SDXL/SD3/FLUX.1 models without any fine-tuning data.
\item \textbf{No learnable parameters}: All projection matrices are identity ($W_Q = W_K = W_V = I$), all geometric operations (exponential map, M\"obius addition, geodesic distance) are closed-form.
\item \textbf{Inference-time optimization}: During the first 5 denoising steps (out of 50 total), composite tokens are iteratively refined via gradient descent on $\mathcal{L}_{\text{sem}} + \lambda_2 \mathcal{L}_{\text{hyp}}$, and attention maps are refined via $\mathcal{L}_{\text{ent}}$.
\item \textbf{Runtime overhead}: Compared to standard SDXL generation ($\sim$8.2\,s), GeoBind adds $\sim$4.6--8.2\,s depending on entity count ($K=2$ to $K=5$), primarily from the additional UNet forward passes required for $\mathcal{L}_{\text{sem}}$ computation.
\item \textbf{Memory}: Peak GPU memory usage is $\sim$18\,GB (SDXL) on A100, dominated by the base model weights. GeoBind's geometric operations add negligible memory overhead ($<$100\,MB).
\end{itemize}

% ============================================================
\section{Code Structure}
\label{sec:code}
% ============================================================

The submitted code contains the following key files:

\paragraph{Core Pipeline.}
\begin{itemize}
\item \texttt{pipe\_geobind.py}: The main GeoBind pipeline, extending \texttt{StableDiffusionXLPipeline} from diffusers. Implements composite token merging (Sec.~3.3), semantic binding loss (Sec.~3.4.1), hyperbolic contrastive loss (Sec.~3.4.2), and entropy-guided attention refinement (Sec.~3.5).
\item \texttt{pipe\_tome.py}: ToMe baseline pipeline (Euclidean-only token merging) for comparison.
\item \texttt{prompt\_utils.py}: SpaCy-based prompt parsing for extracting entity--attribute pairs and aligning with CLIP tokenizer wordpieces.
\item \texttt{utils/hyperbolic\_utils.py}: Hyperbolic geometry operations (exponential/logarithmic maps, M\"obius addition, geodesic distance, Poincar\'e ball projection) and the \texttt{TokenMergerWithAttnHyperspace} module.
\item \texttt{utils/ptp\_utils.py}: Prompt-to-prompt attention utilities for storing and aggregating cross-attention maps.
\end{itemize}

\paragraph{Experiment Scripts.}
\begin{itemize}
\item \texttt{run\_geobind.py}: Generate images on T2I-CompBench attribute binding subsets (color, shape, texture).
\item \texttt{run\_object\_binding\_generate.py}: GOB-Bench image generation for all 200 prompts across 4 difficulty levels.
\item \texttt{run\_object\_binding\_eval.py}: GOB-Bench evaluation with GPT-4o and Qwen2-VL judges, per-difficulty scoring.
\item \texttt{run\_demo.py}: Interactive demo for comparing SDXL, ToMe, and GeoBind on custom prompts.
\end{itemize}

\paragraph{Evaluation Scripts.}
\begin{itemize}
\item \texttt{eval\_image\_reward.py}: ImageReward~\cite{xu2024imagereward} scoring for generated images.
\item \texttt{eval\_clip\_score.py}: CLIP Score (ViT-L/14) evaluation for text-image alignment.
\item \texttt{eval\_gromov\_hyperbolicity.py}: Gromov $\delta$-hyperbolicity analysis of CLIP token embeddings (Table~1 in the main paper), verifying the intrinsic tree-like structure that motivates our hyperbolic approach.
\end{itemize}

\paragraph{Data.}
\begin{itemize}
\item \texttt{data/t2i\_compbench/}: T2I-CompBench prompt files for color, shape, and texture binding evaluation (300 prompts each).
\item \texttt{supplementary/gob\_bench\_prompts.txt}: Complete GOB-Bench prompt list (200 prompts, 4 difficulty levels).
\end{itemize}

\paragraph{BLIP-VQA Evaluation.}
Attribute binding scores (color, texture, shape) are computed using the official BLIP-VQA evaluation script from T2I-CompBench~\cite{huang2023t2icompbench}\footnote{\url{https://github.com/Karine-Huang/T2I-CompBench/tree/main/BLIPvqa_eval}}. Since the BLIP-VQA evaluation requires a different dependency environment (e.g., an older version of \texttt{transformers}), we provide a separate dependency file \texttt{requirements-BLIP-VQA} for setting up this environment. The evaluation workflow is: (1)~generate images using \texttt{run\_geobind.py}; (2)~set up the BLIP-VQA environment; (3)~run the official \texttt{BLIP\_vqa.py} script on the generated images to obtain per-subset BLIP-VQA scores.

GOB-Bench evaluation uses GPT-4o and Qwen2-VL APIs as described in Sec.~4.1 of the main paper.

% ============================================================

\bibliographystyle{ACM-Reference-Format}

\begin{thebibliography}{10}

\bibitem{hu2024tome}
T.~Hu et~al.
\newblock Token Merging for Training-Free Semantic Binding in Text-to-Image Synthesis.
\newblock {\em NeurIPS}, 2024.

\bibitem{rassin2023syngen}
R.~Rassin et~al.
\newblock Linguistic Binding in Diffusion Models: Enhancing Attribute Correspondence through Attention Map Alignment.
\newblock {\em NeurIPS}, 2023.

\bibitem{honnibal2020spacy}
M.~Honnibal et~al.
\newblock spaCy: Industrial-strength Natural Language Processing in Python.
\newblock 2020.

\bibitem{huang2023t2icompbench}
K.~Huang et~al.
\newblock T2I-CompBench: A Comprehensive Benchmark for Open-World Compositional Text-to-Image Generation.
\newblock {\em NeurIPS}, 2023.

\bibitem{hu2024ella}
X.~Hu et~al.
\newblock ELLA: Equip Diffusion Models with LLM for Enhanced Semantic Alignment.
\newblock {\em arXiv:2403.05135}, 2024.

\bibitem{li2024genaibench}
B.~Li et~al.
\newblock GenAI-Bench: Evaluating and Improving Compositional Text-to-Visual Generation.
\newblock {\em arXiv:2406.13743}, 2024.

\bibitem{xu2024imagereward}
J.~Xu et~al.
\newblock ImageReward: Learning and Evaluating Human Preferences for Text-to-Image Generation.
\newblock {\em NeurIPS}, 2024.

\end{thebibliography}

\end{document}
